\documentclass[12pt]{article}

\usepackage{sbc-template}
\usepackage{graphicx,url}
\usepackage[utf8]{inputenc}
\usepackage[brazil]{babel}
\usepackage[latin1]{inputenc}  

\sloppy

\title{Evaluating In-Context Learning Strategies for Text-to-SQL on Local Small Language Models}

\author{Estevan Kuster\inst{1}, Gabriel Ávila\inst{1}, Henrique Trein\inst{1}, Marcelo Stangler\inst{1}}


\address{Instituto de Informática -- Universidade Federal do Rio Grande do Sul
  (UFRGS)\\
  Porto Alegre -- RS -- Brazil
  \email{\{estevan.kuster, gabriel.avila, henrique.trein, marcelo.stangler\}@inf.ufrgs.br}
}

\begin{document} 

\maketitle

\begin{abstract}
  The task of converting natural language questions into SQL queries (Text-to-SQL) has seen significant improvements with Large Language Models (LLMs). However, relying on proprietary, cloud-based models raises concerns regarding data privacy and computational costs. This paper investigates the feasibility of using Small Language Models (SLMs), such as *****, running locally to perform Text-to-SQL tasks on the Spider dataset. We propose an evaluation of In-Context Learning (ICL) strategies, comparing Zero-Shot performance against Few-Shot approaches. Furthermore, we analyze the impact of example selection methods, specifically comparing random selection versus similarity-based retrieval using sentence embeddings. Our methodology aims to demonstrate accessible and private alternatives for effective database querying.
\end{abstract}

\section{Introduction}

Relational databases store a vast amount of the world's data. However, accessing this information requires proficiency in Structured Query Language (SQL), which poses a barrier to non-technical users. The task of Text-to-SQL aims to bridge this gap by automatically translating natural language questions into executable SQL queries. Recent advancements in Natural Language Processing (NLP), driven by Large Language Models (LLMs) such as GPT-4, have achieved state-of-the-art results in code generation benchmarks. Despite their performance, deploying these massive models in real-world enterprise environments is often impractical due to high inference costs, latency, and, crucially, data privacy regulations. This work focuses on evaluating the capabilities of open-source SLMs running locally for the Text-to-SQL task. We utilize the Spider 1.0 dataset, a challenging cross-domain benchmark. Our main contributions are:
\begin{itemize}
    \item An evaluation of SLMs utilizing quantization for local execution on constrained hardware;
    \item A comparative analysis of prompt engineering strategies, specifically Zero-Shot versus Few-Shot learning;
    \item An investigation into the effectiveness of dynamic example selection for Few-Shot prompts, comparing random selection against similarity-based retrieval using sentence embeddings.
\end{itemize}

\section{Methodology} \label{sec:methodology}

Our approach relies on In-Context Learning (ICL), where the model is provided with instructions, the database schema, and optionally, examples of valid Question-SQL pairs within the prompt context window. This method avoids the high computational cost of fine-tuning.

\subsection{Dataset}
We utilize the Spider dataset, a large-scale, complex, cross-domain semantic parsing and text-to-SQL dataset. Spider contains *****.
\subsection{Models and Environment}
To address the constraint of local execution on consumer hardware (e.g., GPUs with limited VRAM), we employ quantized versions of current state-of-the-art SLMs. The models selected for evaluation include ****.

\subsection{Prompt Strategies}
We evaluate two main prompting strategies:
\begin{itemize}
    \item \textbf{Zero-Shot:} The model is provided with the database schema (table names and columns) and the target question, with no prior examples.
    \item \textbf{Few-Shot:} The prompt includes $k$ examples of (Question, SQL) pairs derived from the training set. We define $k=******$ for our experiments to fit within the context window while providing sufficient guidance.
\end{itemize}

\subsection{Example Selection Mechanism}
For the Few-Shot strategy, the selection of examples is critical. We implement two selection mechanisms:
\begin{enumerate}
    \item \textbf{Random Selection:} Examples are blindly sampled from the training set. This serves as our baseline for Few-Shot performance.
    \item \textbf{Similarity-Based Selection:} Inspired by recent works such as DAIL-SQL, we utilize a dense retrieval approach. We *****.

\section{Proposed Evaluation}
The generated SQL queries are evaluated using the official Spider evaluation script. The primary metric is \textit{Execution Accuracy}, which compares the result of executing the predicted SQL query against the result of the ground truth query on the database. 

\bibliographystyle{sbc}
\bibliography{sbc-template}

\end{document}